\chapter{Perfect Objects, Rigid Self-Enrichment, and Canonical Duality}
\label{chap:perfect-rigid-enrichment}
\label{chap:perfect-rigid-poincare}

\section{Purpose and standing conventions}
\label{sec:perfect-rigid-purpose}
\label{sec:perfect-poincare-purpose}

This chapter develops the rigid structures carried by the relative module
and quasi-coherent categories constructed in
\Ref{chap:relative-ag-monoidal-syntax}. Its central object is the full
subcategory
\[
    \Perf(A)
    :=
    \Mod_A(\V)^{\mathrm{dual}}
\]
of dualizable modules over a commutative algebra object
\[
    A\in\CAlg(\V).
\]

The passage from
\[
    \Mod_A(\V)
\]
to
\[
    \Perf(A)
\]
converts the ambient closed symmetric monoidal structure into a rigid
self-enriched structure. In particular, for
\[
    M\in\Perf(A),
    \qquad
    N\in\Mod_A(\V),
\]
the internal Hom is represented by tensoring with the dual:
\[
    \HHom_A(M,N)
    \simeq
    M^\vee\otimes_A N.
\]

The chapter separates and then recombines three layers:
\begin{enumerate}[label=(\roman*)]
    \item the closed self-enrichment of the ambient categories
    \[
        \V
        \qquad\text{and}\qquad
        \Mod_A(\V);
    \]

    \item the rigid stable symmetric monoidal categories
    \[
        \Perf(A)
        \qquad\text{and}\qquad
        \Perf(X),
    \]
    together with their affine and stackwise pullback functoriality;

    \item the canonical symmetric bilinear pairing, the coherent rigid
    duality it represents, enriched dualization, cocomposition, and the
    perfect trace pairing.
\end{enumerate}

The canonical duality is not introduced merely as an objectwise operation.
It is constructed from the symmetric bilinear functor
\[
    \mathbb B_A(M,N)
    :=
    \map_{\Perf(A)}
    \bigl(
        M\otimes_A N,
        A
    \bigr),
\]
whose symmetry is induced by the braiding of
\[
    \Perf(A).
\]
Its right representing functor is dualization,
\[
    D_A
    :=
    (-)^\vee
    \colon
    \Perf(A)^{\op}
    \xrightarrow{\;\simeq\;}
    \Perf(A).
\]
The coherent symmetry of
\[
    \mathbb B_A
\]
supplies
\[
    D_A
\]
with its full coherent duality structure, including biduality and all
higher
\[
    C_2
\]
-coherences.

The categorical, mapping-space, coherence, universe, and sheaf conventions
of
\Ref{chap:relative-ag-monoidal-syntax}
remain in force.

\begin{convention}[Higher coherence in the rigid layer]
\label{conv:rigid-higher-coherence}
In addition to
\Ref{conv:higher-coherence},
duality, biduality, enrichment, enriched dualization, cocomposition, trace,
and cyclicity data are always understood together with all required
coherent higher homotopies.

Whenever such data are suppressed from the notation, their existence is
supplied by a rigid, enriched, or symmetric-bilinear construction, by a
universal property with contractible choice space, or by an explicitly
cited coherence theorem.
\end{convention}

\begin{convention}[Perfect categorical tensor products]
\label{conv:perfect-categorical-tensor-product}
The notation
\[
    \Cat_\infty^{\mathrm{perf}}
\]
denotes the $\infty$-category of small stable idempotent-complete
$\infty$-categories and exact functors. It is equipped with the symmetric
monoidal tensor product characterized by the natural equivalence
\[
    \Fun^{\mathrm{ex}}
    \bigl(
        \mathcal C\otimes\mathcal D,
        \mathcal E
    \bigr)
    \simeq
    \Fun^{\mathrm{biex}}
    \bigl(
        \mathcal C\times\mathcal D,
        \mathcal E
    \bigr),
\]
where the right-hand side denotes functors exact separately in each
variable.

Accordingly,
\[
    \CAlg(\Cat_\infty^{\mathrm{perf}})
\]
is formed using this tensor product. Its objects are small stable
idempotent-complete symmetric monoidal $\infty$-categories whose tensor
products are exact separately in each variable, and its morphisms are
exact strong symmetric monoidal functors.

Let
\[
    \omega_{\mathrm{perf}}
    \colon
    \CAlg(\Cat_\infty^{\mathrm{perf}})
    \longrightarrow
    \CAlg(\Cat_\infty)
\]
denote the structural forgetful functor. It sends a small stable
idempotent-complete symmetric monoidal $\infty$-category with biexact
tensor product to its underlying symmetric monoidal $\infty$-category,
and it sends an exact strong symmetric monoidal functor to its underlying
strong symmetric monoidal functor.

Thus
\[
    \omega_{\mathrm{perf}}
\]
forgets stability, idempotent completeness, and exactness while retaining
the underlying symmetric monoidal structures and strong symmetric
monoidal functors. It is not being defined by applying
\[
    \CAlg(-)
\]
to a symmetric monoidal inclusion between
\[
    \Cat_\infty^{\mathrm{perf}}
\]
with its perfect categorical tensor product and
\[
    \Cat_\infty
\]
with its Cartesian symmetric monoidal structure.

All occurrences of
\[
    \Cat_\infty^{\mathrm{perf}}
\]
are formed in a universe sufficiently large that every perfect category
under consideration is small in that universe.
\end{convention}

\section{Closed self-enrichment of the ambient module categories}
\label{sec:closed-self-enrichment-before-perfect-restriction}

The ambient closed symmetric monoidal structures already carry canonical
self-enrichments. The perfect layer is obtained by restricting those
enrichments to dualizable objects.

\begin{proposition}[Canonical self-enrichment]
\label{prop:ambient-self-enrichment}
The closed symmetric monoidal structure on
\[
    \V
\]
determines a canonical
\[
    \V
\]
-enriched $\infty$-category whose objects are the objects of
\[
    \V
\]
and whose enriched Hom objects are
\[
    \underline{\V}(X,Y)
    :=
    \HHom_{\V}(X,Y).
\]

For every finite string
\[
    X_0,X_1,\ldots,X_n
    \in
    \V
\]
with
\[
    n\geq1,
\]
there is a canonical $n$-ary composition morphism
\[
\begin{aligned}
    c_{X_0,\ldots,X_n}
    \colon{}&
    \HHom_{\V}(X_{n-1},X_n)
    \otimes
    \cdots
    \otimes
    \HHom_{\V}(X_0,X_1)
    \\
    &\longrightarrow
    \HHom_{\V}(X_0,X_n).
\end{aligned}
\]
For
\[
    n=0,
\]
the corresponding nullary operation is the enriched identity
\[
    \mathbf 1_{\V}
    \longrightarrow
    \HHom_{\V}(X_0,X_0).
\]
These operations satisfy all associative and unital higher coherences.
\end{proposition}

\begin{proof}
By
\Ref{prop:ambient-closedness},
the symmetric monoidal $\infty$-category
\[
    \V
\]
is closed. The canonical self-enrichment construction for a closed
symmetric monoidal $\infty$-category therefore supplies a
\[
    \V
\]
-enriched $\infty$-category whose enriched Hom objects are
\[
    \underline{\V}(X,Y)
    =
    \HHom_{\V}(X,Y).
\]

For
\[
    n\geq1,
\]
consider the iterated evaluation morphism
\[
\begin{aligned}
    &
    \HHom_{\V}(X_{n-1},X_n)
    \otimes
    \cdots
    \otimes
    \HHom_{\V}(X_0,X_1)
    \otimes
    X_0
    \\
    &\qquad\longrightarrow
    \HHom_{\V}(X_{n-1},X_n)
    \otimes
    \cdots
    \otimes
    X_1
    \\
    &\qquad\longrightarrow
    \cdots
    \longrightarrow
    X_n.
\end{aligned}
\]
Adjunction with respect to
\[
    -\otimes X_0
    \dashv
    \HHom_{\V}(X_0,-)
\]
gives the composition morphism
\[
\begin{aligned}
    c_{X_0,\ldots,X_n}
    \colon{}&
    \HHom_{\V}(X_{n-1},X_n)
    \otimes
    \cdots
    \otimes
    \HHom_{\V}(X_0,X_1)
    \\
    &\longrightarrow
    \HHom_{\V}(X_0,X_n).
\end{aligned}
\]

For
\[
    n=0,
\]
the enriched identity
\[
    \mathbf 1_{\V}
    \longrightarrow
    \HHom_{\V}(X_0,X_0)
\]
is adjoint to the unit equivalence
\[
    \mathbf 1_{\V}\otimes X_0
    \simeq
    X_0.
\]

Let
\[
    0=i_0<i_1<\cdots<i_r=n
\]
be a subdivision of the string
\[
    X_0,\ldots,X_n.
\]
One may first compose inside the blocks
\[
    X_{i_{j-1}},\ldots,X_{i_j}
\]
and then compose the resulting enriched morphisms, or apply
\[
    c_{X_0,\ldots,X_n}
\]
directly. After tensoring either resulting map with
\[
    X_0,
\]
both composites are canonically identified with the same iterated
evaluation morphism
\[
\begin{aligned}
    &
    \HHom_{\V}(X_{n-1},X_n)
    \otimes
    \cdots
    \otimes
    \HHom_{\V}(X_0,X_1)
    \otimes
    X_0
    \longrightarrow
    X_n.
\end{aligned}
\]
Adjunction identifies the corresponding composition maps.

For nested subdivisions, the comparison homotopies are obtained by taking
adjoints of the associativity and unit coherences of iterated evaluation.
The canonical self-enrichment construction supplies all higher
compatibilities. See
\cite{GepnerHaugsengEnrichedInfinityCategories,LurieHigherAlgebra}.
\end{proof}

\begin{remark}
\label{rem:closed-not-rigid}
The existence of the internal Hom does not imply that every object of
\[
    \V
\]
is dualizable. The equivalence
\[
    \HHom_{\V}(X,Y)
    \simeq
    X^\vee\otimes Y
\]
holds when
\[
    X
\]
is dualizable.
\end{remark}

\begin{proposition}[Self-enrichment of the module category]
\label{prop:module-self-enrichment}
The category
\[
    \Mod_A(\V)
\]
is canonically enriched over itself, with enriched Hom objects
\[
    \underline{\Mod_A(\V)}(M,N)
    =
    \HHom_A(M,N).
\]
\end{proposition}

\begin{proof}
Apply
\Ref{prop:ambient-self-enrichment}
to the presentably stable symmetric monoidal category
\[
    \Mod_A(\V).
\]
\end{proof}

\section{Perfect modules and tensor-dual internal Homs}
\label{sec:quasi-coherent-perfect}


\subsection{Dualizable and perfect modules}

\begin{definition}[Dualizable module]
\label{def:dualizable-module}
An
\[
    A
\]
-module
\[
    M
\]
is \emph{dualizable} if there exists an
\[
    A
\]
-module
\[
    M^\vee
\]
and morphisms
\[
    \coev_M
    \colon
    A
    \longrightarrow
    M\otimes_A M^\vee,
\]
\[
    \ev_M
    \colon
    M^\vee\otimes_A M
    \longrightarrow
    A
\]
satisfying the two triangle identities together with their coherent
higher homotopies.
\end{definition}

\begin{definition}[Perfect module]
\label{def:perfect-module}
The full subcategory of perfect
\[
    A
\]
-modules is
\[
    \Perf(A)
    :=
    \Mod_A(\V)^{\mathrm{dual}}.
\]
\end{definition}

\begin{remark}[Perfect versus compact]
\label{rem:perfect-versus-compact}
Perfectness is defined here by dualizability.

For an $\EE_\infty$-ring spectrum
\[
    A,
\]
an object of
\[
    \Mod_A(\Sp)
\]
is dualizable if and only if it is compact. Thus the present definition
recovers the usual category of perfect $A$-modules.

For module categories internal to a general presentably stable symmetric
monoidal category
\[
    \V,
\]
compactness and dualizability need not coincide. No such identification is
used in this chapter.
\end{remark}


\subsection{Closure properties of perfect modules}

\begin{lemma}[Closure properties of perfect modules]
\label{lem:perfect-closed-tensor-dual}
The category
\[
    \Perf(A)
\]
is closed under equivalences, finite tensor products, duals, all shifts,
cofibres, and retracts.
\end{lemma}

\begin{proof}
Closure under equivalences is immediate.

If
\[
    M
    \qquad\text{and}\qquad
    N
\]
are dualizable, then
\[
    M\otimes_A N
\]
is dualizable with
\[
    (M\otimes_A N)^\vee
    \simeq
    N^\vee\otimes_A M^\vee.
\]
Evaluation and coevaluation are obtained by tensoring the duality data for
\[
    M
    \qquad\text{and}\qquad
    N
\]
and contracting the corresponding dual pairs.

The dual of a dualizable object is dualizable, and the canonical
biduality map gives
\[
    (M^\vee)^\vee
    \simeq
    M.
\]

The zero object is dualizable, with dual the zero object.

We next prove closure under all shifts. Since tensoring is exact, there are
canonical equivalences
\[
    (\Sigma^rX)\otimes_A(\Sigma^sY)
    \simeq
    \Sigma^{r+s}(X\otimes_A Y)
\]
for all
\[
    r,s\in\mathbb Z.
\]
In particular,
\[
    (\Sigma^nA)\otimes_A(\Sigma^{-n}A)
    \simeq
    A
\]
and
\[
    (\Sigma^{-n}A)\otimes_A(\Sigma^nA)
    \simeq
    A.
\]
Thus
\[
    \Sigma^nA
\]
is tensor-invertible, with tensor inverse
\[
    \Sigma^{-n}A,
\]
and is therefore dualizable. For every dualizable
\[
    M,
\]
one has
\[
    \Sigma^nM
    \simeq
    (\Sigma^nA)\otimes_A M.
\]
Hence
\[
    \Sigma^nM
\]
is dualizable for every
\[
    n\in\mathbb Z.
\]
Moreover,
\[
\begin{aligned}
    (\Sigma^nM)^\vee
    &\simeq
    \bigl(
        (\Sigma^nA)\otimes_A M
    \bigr)^\vee
    \\
    &\simeq
    M^\vee\otimes_A(\Sigma^nA)^\vee
    \\
    &\simeq
    M^\vee\otimes_A\Sigma^{-n}A
    \\
    &\simeq
    \Sigma^{-n}M^\vee.
\end{aligned}
\]

We next prove closure under cofibres. Let
\[
    M
    \xrightarrow{u}
    N
    \longrightarrow
    P
\]
be a cofiber sequence with
\[
    M
    \qquad\text{and}\qquad
    N
\]
dualizable. Dualizing
\[
    u
\]
gives a morphism
\[
    u^\vee
    \colon
    N^\vee
    \longrightarrow
    M^\vee.
\]
Define
\[
    P^\vee
    :=
    \fib(u^\vee).
\]

For every
\[
    X\in\Mod_A(\V),
\]
exactness of tensoring gives
\[
\begin{aligned}
    P^\vee\otimes_A X
    &\simeq
    \fib
    \left(
        N^\vee\otimes_A X
        \longrightarrow
        M^\vee\otimes_A X
    \right)
    \\
    &\simeq
    \fib
    \left(
        \HHom_A(N,X)
        \longrightarrow
        \HHom_A(M,X)
    \right).
\end{aligned}
\]

We identify the last fibre with
\[
    \HHom_A(P,X).
\]
For every
\[
    L\in\Mod_A(\V),
\]
one has
\[
\begin{aligned}
    &
    \Map_{\Mod_A(\V)}
    \left(
        L,
        \fib
        \left(
            \HHom_A(N,X)
            \longrightarrow
            \HHom_A(M,X)
        \right)
    \right)
    \\
    &\qquad\simeq
    \fib
    \left(
        \Map_{\Mod_A(\V)}
        \bigl(
            L\otimes_A N,
            X
        \bigr)
        \longrightarrow
        \Map_{\Mod_A(\V)}
        \bigl(
            L\otimes_A M,
            X
        \bigr)
    \right)
    \\
    &\qquad\simeq
    \Map_{\Mod_A(\V)}
    \bigl(
        L\otimes_A P,
        X
    \bigr)
    \\
    &\qquad\simeq
    \Map_{\Mod_A(\V)}
    \bigl(
        L,
        \HHom_A(P,X)
    \bigr).
\end{aligned}
\]
The middle equivalence uses exactness of
\[
    L\otimes_A-.
\]
The Yoneda lemma therefore gives a natural equivalence
\[
    P^\vee\otimes_A X
    \xrightarrow{\;\simeq\;}
    \HHom_A(P,X).
\]

Consequently,
\[
    -\otimes_A P
    \dashv
    -\otimes_A P^\vee.
\]
The unit and counit of this adjunction, evaluated at
\[
    A,
\]
are morphisms
\[
    \coev_P
    \colon
    A
    \longrightarrow
    P\otimes_A P^\vee
\]
and
\[
    \ev_P
    \colon
    P^\vee\otimes_A P
    \longrightarrow
    A.
\]
Their triangle identities are the unit--counit identities of the
adjunction. Hence
\[
    P
\]
is dualizable with dual
\[
    P^\vee.
\]

Finally, suppose that
\[
    M
    \xrightarrow{i}
    N
    \xrightarrow{r}
    M
\]
exhibits
\[
    M
\]
as a retract of a dualizable object
\[
    N,
\]
so that
\[
    r i
    \simeq
    \id_M.
\]
Let
\[
    e
    :=
    i r
    \colon
    N
    \longrightarrow
    N.
\]
The dual idempotent
\[
    e^\vee
    \colon
    N^\vee
    \longrightarrow
    N^\vee
\]
splits because
\[
    \Mod_A(\V)
\]
is idempotent complete. Choose a splitting
\[
    M'
    \xrightarrow{j}
    N^\vee
    \xrightarrow{s}
    M'
\]
with
\[
    s j
    \simeq
    \id_{M'},
    \qquad
    j s
    \simeq
    e^\vee.
\]

Define
\[
    \ev_M
    \colon
    M'\otimes_A M
    \xrightarrow{j\otimes i}
    N^\vee\otimes_A N
    \xrightarrow{\ev_N}
    A
\]
and
\[
    \coev_M
    \colon
    A
    \xrightarrow{\coev_N}
    N\otimes_A N^\vee
    \xrightarrow{r\otimes s}
    M\otimes_A M'.
\]

The idempotents
\[
    e
    \qquad\text{and}\qquad
    e^\vee
\]
are mates under the duality of
\[
    N.
\]
Consequently,
\[
    \ev_N
    \circ
    \bigl(
        e^\vee\otimes\id_N
    \bigr)
    \simeq
    \ev_N
    \circ
    \bigl(
        \id_{N^\vee}\otimes e
    \bigr)
\]
and
\[
    \bigl(
        e\otimes\id_{N^\vee}
    \bigr)
    \circ
    \coev_N
    \simeq
    \bigl(
        \id_N\otimes e^\vee
    \bigr)
    \circ
    \coev_N.
\]
Using
\[
    r i\simeq\id_M,
    \qquad
    i r\simeq e,
    \qquad
    s j\simeq\id_{M'},
    \qquad
    j s\simeq e^\vee,
\]
the two triangle composites for
\[
    M
    \qquad\text{and}\qquad
    M'
\]
reduce to the triangle identities for
\[
    N
    \qquad\text{and}\qquad
    N^\vee.
\]
Thus
\[
    M'
\]
is a dual of
\[
    M.
\]
\end{proof}


\subsection{The stable rigid category of perfect modules}

\begin{proposition}[Stability of the perfect category]
\label{prop:perfect-stable}
The category
\[
    \Perf(A)
\]
is a stable idempotent-complete symmetric monoidal $\infty$-category.
Its tensor product is exact separately in each variable.
\end{proposition}

\begin{proof}
The tensor unit
\[
    A
\]
is dualizable. The zero object is dualizable, with dual the zero object.
By
\Ref{lem:perfect-closed-tensor-dual},
the full subcategory
\[
    \Perf(A)
\]
is closed under all shifts, cofibres, tensor products, and retracts. It is
therefore stable and idempotent complete.

Its tensor product is the restriction of the tensor product on
\[
    \Mod_A(\V).
\]
Since tensoring in
\[
    \Mod_A(\V)
\]
is exact separately in each variable, its restriction to
\[
    \Perf(A)
\]
is exact separately in each variable.
\end{proof}


\subsection{Tensor-dual realization of internal Homs}

\begin{theorem}[Tensor-dual realization of the internal Hom]
\label{thm:tensor-dual-internal-hom}
Let
\[
    M\in\Perf(A),
    \qquad
    N\in\Mod_A(\V).
\]
There is a canonical natural equivalence
\[
    M^\vee\otimes_A N
    \xrightarrow{\;\simeq\;}
    \HHom_A(M,N).
\]
Equivalently, there is a canonical equivalence of functors
\[
    M^\vee\otimes_A-
    \xrightarrow{\;\simeq\;}
    \HHom_A(M,-).
\]
\end{theorem}

\begin{proof}
The duality of
\[
    M
\]
gives an adjunction
\[
    -\otimes_A M
    \dashv
    -\otimes_A M^\vee.
\]
The closed structure on
\[
    \Mod_A(\V)
\]
gives an adjunction
\[
    -\otimes_A M
    \dashv
    \HHom_A(M,-).
\]
Right adjoints to a fixed functor are unique through a contractible space
of choices. Therefore there is a canonical natural equivalence
\[
    M^\vee\otimes_A-
    \xrightarrow{\;\simeq\;}
    \HHom_A(M,-).
\]
Evaluating at
\[
    N
\]
gives the stated equivalence.
\end{proof}

\begin{corollary}
\label{cor:perfect-internal-hom}
If
\[
    M,N\in\Perf(A),
\]
then
\[
    \HHom_A(M,N)
    \in
    \Perf(A).
\]
\end{corollary}

\begin{proof}
By
\Ref{thm:tensor-dual-internal-hom},
\[
    \HHom_A(M,N)
    \simeq
    M^\vee\otimes_A N.
\]
The right-hand side is dualizable by
\Ref{lem:perfect-closed-tensor-dual}.
\end{proof}

\section{Indexed perfect categories and perfect objects on relative stacks}
\label{sec:indexed-perfect-stacks}


\subsection{Affine base change and indexed perfect categories}

\begin{proposition}[Base change of perfect modules]
\label{prop:base-change-perfect-modules}
Let
\[
    A
    \longrightarrow
    B
\]
be a morphism in
\[
    \CAlg(\V).
\]
If
\[
    M\in\Perf(A),
\]
then
\[
    M_B
    :=
    M\otimes_A B
\]
belongs to
\[
    \Perf(B),
\]
and there is a canonical equivalence
\[
    (M_B)^\vee
    \simeq
    M^\vee\otimes_A B.
\]
\end{proposition}

\begin{proof}
Extension of scalars is strong symmetric monoidal by
\Ref{prop:monoidality-base-change}.
A strong symmetric monoidal functor carries evaluation, coevaluation, and
the triangle identities to evaluation, coevaluation, and the triangle
identities. It therefore preserves dualizable objects and their duals.
\end{proof}

\begin{proposition}[The indexed perfect category]
\label{prop:indexed-perfect-category}
The indexed module category restricts to a functor
\[
    \mathbf{Perf}_{\V}
    \colon
    \CAlg(\V)
    \longrightarrow
    \CAlg(\Cat_\infty^{\mathrm{perf}})
\]
given by
\[
    A
    \longmapsto
    \Perf(A),
\]
and
\[
    (A\to B)
    \longmapsto
    (-\otimes_A B).
\]
\end{proposition}

\begin{proof}
By
\Ref{prop:base-change-perfect-modules},
extension of scalars preserves dualizable objects and their duality data.
Hence every morphism
\[
    A
    \longrightarrow
    B
\]
induces a strong symmetric monoidal functor
\[
    -\otimes_A B
    \colon
    \Perf(A)
    \longrightarrow
    \Perf(B).
\]

Extension of scalars preserves all small colimits on the ambient module
categories. Since both source and target are stable, it is exact and
therefore preserves zero objects, fibre sequences, and cofiber sequences.
Its restriction to perfect objects is consequently exact.

The coherence of composition and units is inherited from the indexed
module category of
\Ref{prop:module-indexed-category}.
Thus the indexed module functor restricts to
\[
    \mathbf{Perf}_{\V}
    \colon
    \CAlg(\V)
    \longrightarrow
    \CAlg(\Cat_\infty^{\mathrm{perf}}).
\]
The target is interpreted using
\Ref{conv:perfect-categorical-tensor-product}.
\end{proof}

\subsection{Perfect objects on relative stacks}

\begin{definition}[Perfect quasi-coherent object]
\label{def:perfect-qcoh-object}
For
\[
    X\in\PSh(\Aff),
\]
define
\[
    \Perf(X)
    :=
    \QCoh(X)^{\mathrm{dual}}.
\]
\end{definition}

\begin{lemma}[Duality data commute with symmetric monoidal limits]
\label{lem:duality-data-commute-with-limits}
Let
\[
    F
    \colon
    I
    \longrightarrow
    \CAlg(\Cat_\infty),
    \qquad
    i
    \longmapsto
    \mathcal C_i,
\]
be a small diagram of symmetric monoidal $\infty$-categories and strong
symmetric monoidal functors. Let
\[
    \mathcal C
    :=
    \lim_{i\in I}\mathcal C_i
\]
be its limit in
\[
    \CAlg(\Cat_\infty).
\]

For
\[
    M=(M_i)_{i\in I}\in\mathcal C,
\]
let
\[
    \Dual_{\mathcal C}(M)
\]
denote the space of choices of a dual of
\[
    M,
\]
including evaluation, coevaluation, the triangle homotopies, and all
coherent higher data.

Restriction to the factors induces a natural equivalence
\[
    \Dual_{\mathcal C}(M)
    \xrightarrow{\;\simeq\;}
    \lim_{i\in I}
    \Dual_{\mathcal C_i}(M_i).
\]

Consequently, an object
\[
    M\in\mathcal C
\]
is dualizable if and only if every component
\[
    M_i\in\mathcal C_i
\]
is dualizable.

Moreover, restriction induces an equivalence of symmetric monoidal
$\infty$-categories
\[
    \mathcal C^{\mathrm{dual}}
    \xrightarrow{\;\simeq\;}
    \lim_{i\in I}
    \mathcal C_i^{\mathrm{dual}}
\]
in
\[
    \CAlg(\Cat_\infty).
\]
\end{lemma}

\begin{proof}
Let
\[
    \operatorname{DualPair}^{\otimes}
\]
denote the walking dual pair: the symmetric monoidal $\infty$-category
freely generated by two objects
\[
    x,
    \qquad
    x^\vee,
\]
together with evaluation and coevaluation morphisms
\[
    x^\vee\otimes x
    \longrightarrow
    \mathbf 1,
\]
\[
    \mathbf 1
    \longrightarrow
    x\otimes x^\vee,
\]
and the coherent triangle identities.

For a symmetric monoidal $\infty$-category
\[
    \mathcal D,
\]
let
\[
    \operatorname{DP}(\mathcal D)
    :=
    \Map_{\CAlg(\Cat_\infty)}
    \bigl(
        \operatorname{DualPair}^{\otimes},
        \mathcal D
    \bigr)
\]
be its space of dual pairs.

Evaluation at the generating object
\[
    x
\]
defines a morphism
\[
    e_x
    \colon
    \operatorname{DP}(\mathcal D)
    \longrightarrow
    \mathcal D^\simeq.
\]
For
\[
    Y\in\mathcal D,
\]
the space
\[
    \Dual_{\mathcal D}(Y)
\]
is the homotopy fibre of
\[
    e_x
\]
over
\[
    Y\in\mathcal D^\simeq.
\]

Since mapping out of a fixed object preserves limits in the target, the
limit presentation
\[
    \mathcal C
    \simeq
    \lim_{i\in I}\mathcal C_i
\]
induces an equivalence
\[
\begin{aligned}
    \operatorname{DP}(\mathcal C)
    &=
    \Map_{\CAlg(\Cat_\infty)}
    \bigl(
        \operatorname{DualPair}^{\otimes},
        \mathcal C
    \bigr)
    \\
    &\simeq
    \Map_{\CAlg(\Cat_\infty)}
    \left(
        \operatorname{DualPair}^{\otimes},
        \lim_{i\in I}\mathcal C_i
    \right)
    \\
    &\simeq
    \lim_{i\in I}
    \Map_{\CAlg(\Cat_\infty)}
    \bigl(
        \operatorname{DualPair}^{\otimes},
        \mathcal C_i
    \bigr)
    \\
    &=
    \lim_{i\in I}
    \operatorname{DP}(\mathcal C_i).
\end{aligned}
\]

The maximal-subgroupoid functor
\[
    (-)^\simeq
    \colon
    \Cat_\infty
    \longrightarrow
    \mathcal S
\]
is a right adjoint and therefore preserves limits. Hence
\[
    \mathcal C^\simeq
    \simeq
    \lim_{i\in I}
    \mathcal C_i^\simeq.
\]

The evaluation morphism
\[
    e_x
    \colon
    \operatorname{DP}(\mathcal C)
    \longrightarrow
    \mathcal C^\simeq
\]
is, under these equivalences, the limit of the evaluation morphisms
\[
    e_{x,i}
    \colon
    \operatorname{DP}(\mathcal C_i)
    \longrightarrow
    \mathcal C_i^\simeq.
\]
Taking the homotopy fibre over the compatible object
\[
    M=(M_i)_{i\in I}
\]
and using that homotopy fibres are limits in
\[
    \mathcal S
\]
gives the natural equivalence
\[
    \Dual_{\mathcal C}(M)
    \xrightarrow{\;\simeq\;}
    \lim_{i\in I}
    \Dual_{\mathcal C_i}(M_i).
\]

If
\[
    M
\]
is dualizable in
\[
    \mathcal C,
\]
then every
\[
    M_i
\]
is dualizable because each evaluation functor
\[
    \ev_i
    \colon
    \mathcal C
    \longrightarrow
    \mathcal C_i
\]
is strong symmetric monoidal.

Conversely, suppose that every
\[
    M_i
\]
is dualizable. For a fixed dualizable object, the space of complete
duality data is contractible. Thus every space
\[
    \Dual_{\mathcal C_i}(M_i)
\]
is contractible. The diagram
\[
    i
    \longmapsto
    \Dual_{\mathcal C_i}(M_i)
\]
is therefore objectwise equivalent to the constant diagram at the terminal
space. Its limit is contractible and in particular nonempty. The preceding
equivalence then shows that
\[
    \Dual_{\mathcal C}(M)
\]
is nonempty. Hence
\[
    M
\]
is dualizable.

We now compare the full subcategories of dualizable objects.

Restriction along the evaluation functors defines a canonical strong
symmetric monoidal functor
\[
    r
    \colon
    \mathcal C^{\mathrm{dual}}
    \longrightarrow
    \lim_{i\in I}
    \mathcal C_i^{\mathrm{dual}}.
\]

Conversely, the inclusions
\[
    \mathcal C_i^{\mathrm{dual}}
    \hookrightarrow
    \mathcal C_i
\]
induce a canonical strong symmetric monoidal functor
\[
    j
    \colon
    \lim_{i\in I}
    \mathcal C_i^{\mathrm{dual}}
    \longrightarrow
    \lim_{i\in I}
    \mathcal C_i
    \simeq
    \mathcal C.
\]
Every object in the source has dualizable components. By the componentwise
criterion proved above, its image in
\[
    \mathcal C
\]
is dualizable. Therefore
\[
    j
\]
factors canonically through a strong symmetric monoidal functor
\[
    \bar j
    \colon
    \lim_{i\in I}
    \mathcal C_i^{\mathrm{dual}}
    \longrightarrow
    \mathcal C^{\mathrm{dual}}.
\]

The composite
\[
    r\circ\bar j
\]
is canonically equivalent to the identity because its composite with every
evaluation functor
\[
    \lim_{i\in I}
    \mathcal C_i^{\mathrm{dual}}
    \longrightarrow
    \mathcal C_i^{\mathrm{dual}}
\]
is the corresponding evaluation functor.

Similarly,
\[
    \bar j\circ r
\]
is canonically equivalent to the identity on
\[
    \mathcal C^{\mathrm{dual}},
\]
because the inclusion
\[
    \mathcal C^{\mathrm{dual}}
    \hookrightarrow
    \mathcal C
\]
is fully faithful and the two functors become canonically equivalent after
composition with every evaluation
\[
    \mathcal C
    \longrightarrow
    \mathcal C_i.
\]
The universal property of the limit supplies the resulting coherent
equivalences of functors.

Thus
\[
    r
    \colon
    \mathcal C^{\mathrm{dual}}
    \xrightarrow{\;\simeq\;}
    \lim_{i\in I}
    \mathcal C_i^{\mathrm{dual}}
\]
is an equivalence in
\[
    \CAlg(\Cat_\infty).
\]
\end{proof}


\subsection{Affine-point reconstruction}

\begin{proposition}[Perfect objects are computed from affine points]
\label{prop:perfect-limit-description}
Let
\[
    X\in\PSh(\Aff).
\]
The symmetric monoidal limit presentation
\[
    \QCoh(X)
    \simeq
    \lim_{
        (U\to X)\in(\Aff_{/X})^{\op}
    }
    \QCoh(U)
\]
induces an equivalence
\[
    \Perf(X)
    \xrightarrow{\;\simeq\;}
    \lim_{
        (U\to X)\in(\Aff_{/X})^{\op}
    }
    \Perf(U)
\]
in
\[
    \CAlg(\Cat_\infty).
\]

More precisely, for every
\[
    M\in\QCoh(X),
\]
there is a natural equivalence
\[
    \Dual_{\QCoh(X)}(M)
    \simeq
    \lim_{
        (U\to X)\in(\Aff_{/X})^{\op}
    }
    \Dual_{\QCoh(U)}(M|_U).
\]
In particular,
\[
    M
\]
is perfect if and only if every restriction
\[
    M|_U
\]
is perfect.
\end{proposition}

\begin{proof}
The displayed description of
\[
    \QCoh(X)
\]
is a limit in
\[
    \CAlg(\Cat_\infty),
\]
whose transition functors are strong symmetric monoidal
extension-of-scalars functors. Applying
\Ref{lem:duality-data-commute-with-limits}
gives
\[
    \Dual_{\QCoh(X)}(M)
    \simeq
    \lim_{
        (U\to X)\in(\Aff_{/X})^{\op}
    }
    \Dual_{\QCoh(U)}(M|_U).
\]
It follows that
\[
    M
\]
is dualizable in
\[
    \QCoh(X)
\]
if and only if every restriction
\[
    M|_U
\]
is dualizable in
\[
    \QCoh(U).
\]

The tensor products and transition functors on the full subcategories of
dualizable objects are the restrictions of the corresponding structures
on the ambient categories. Hence
\Ref{lem:duality-data-commute-with-limits}
also supplies the asserted equivalence in
\[
    \CAlg(\Cat_\infty).
\]
\end{proof}

\begin{remark}[Affine points versus atlas descent]
\label{rem:perfect-affine-points-versus-descent}
The equivalence of
\Ref{prop:perfect-limit-description}
is indexed by the full category of affine points of
\[
    X.
\]
It does not by itself imply that
\[
    \Perf(X)
\]
can be computed from a selected
\[
    \tau
\]
-atlas or from its \v{C}ech nerve. Such a reduction requires the
corresponding descent theorem for
\[
    \QCoh.
\]
\end{remark}


\subsection{Stability, rigidity, and pullback functoriality}

\begin{proposition}[Stability and monoidality of perfect objects]
\label{prop:perfect-stable-monoidal}
For every
\[
    X\in\PSh(\Aff),
\]
the $\infty$-category
\[
    \Perf(X)
\]
is stable and idempotent complete. Its tensor product is exact separately
in each variable, and the inclusion
\[
    \Perf(X)
    \hookrightarrow
    \QCoh(X)
\]
is exact and strong symmetric monoidal.

Consequently,
\[
    \Perf(X)
\]
is a rigid stable symmetric monoidal $\infty$-category. In particular, it
is canonically enriched over itself, with enriched Hom object
\[
    \underline{\Perf(X)}(M,N)
    \simeq
    M^\vee\otimes N.
\]

Moreover, the equivalence of
\Ref{prop:perfect-limit-description}
upgrades canonically to an equivalence
\[
    \Perf(X)
    \xrightarrow{\;\simeq\;}
    \lim_{
        (U\to X)\in(\Aff_{/X})^{\op}
    }
    \Perf(U)
\]
in
\[
    \CAlg(\Cat_\infty^{\mathrm{perf}}).
\]
\end{proposition}

\begin{proof}
By
\Ref{prop:perfect-limit-description},
there is an equivalence of symmetric monoidal $\infty$-categories
\[
    \Perf(X)
    \simeq
    \lim_{
        (U\to X)\in(\Aff_{/X})^{\op}
    }
    \Perf(U)
\]
in
\[
    \CAlg(\Cat_\infty).
\]

After forgetting the symmetric monoidal structures, the right-hand side
is a limit in
\[
    \Cat_\infty.
\]
Every category
\[
    \Perf(U)
\]
is stable, and every transition functor is exact. Hence zero objects,
fibres, and cofibres in the limit are computed componentwise.

More explicitly, let
\[
    M\longrightarrow N\longrightarrow P
\]
be a diagram in the limit. It is a cofiber sequence if and only if its
evaluation at every affine point is a cofiber sequence. Since every
\[
    \Perf(U)
\]
is stable, the evaluated sequence is also a fibre sequence. Limits and
equivalences in the limit category are detected componentwise, so the
original sequence is a fibre sequence. Thus the limit category is stable.

We next verify idempotent completeness. Let
\[
    e
    \colon
    M
    \longrightarrow
    M
\]
be an idempotent in the limit. Every component
\[
    e_U
    \colon
    M_U
    \longrightarrow
    M_U
\]
splits because
\[
    \Perf(U)
\]
is idempotent complete.

The space of splittings of an idempotent is obtained as a fibre of the
space of functors from the walking retraction
\[
    M
    \xrightarrow{i}
    N
    \xrightarrow{r}
    M,
    \qquad
    r i\simeq\id_M.
\]
Mapping out of this fixed finite category preserves limits. Consequently,
the space of splittings of
\[
    e
\]
is the limit of the spaces of splittings of the
\[
    e_U.
\]
Each of the latter spaces is contractible. Their limit is therefore
contractible and in particular nonempty. Hence
\[
    e
\]
splits, and
\[
    \Perf(X)
\]
is idempotent complete.

The tensor product is computed componentwise:
\[
    (M\otimes N)|_U
    \simeq
    M|_U\otimes N|_U.
\]
For fixed
\[
    M,
\]
let
\[
    N'\longrightarrow N\longrightarrow N''
\]
be a cofiber sequence in
\[
    \Perf(X).
\]
After evaluation at every affine point,
\[
    M|_U\otimes N'|_U
    \longrightarrow
    M|_U\otimes N|_U
    \longrightarrow
    M|_U\otimes N''|_U
\]
is a cofiber sequence because tensoring in
\[
    \Perf(U)
\]
is exact. Therefore
\[
    M\otimes N'
    \longrightarrow
    M\otimes N
    \longrightarrow
    M\otimes N''
\]
is a cofiber sequence in
\[
    \Perf(X).
\]
The same argument applies in the other variable.

The inclusion
\[
    \Perf(X)
    \hookrightarrow
    \QCoh(X)
\]
preserves zero objects and cofiber sequences because these are computed in
the ambient category and detected componentwise. It is therefore exact.

The monoidal unit of
\[
    \QCoh(X)
\]
is dualizable, and dualizable objects are closed under tensor products.
Hence
\[
    \Perf(X)
    \hookrightarrow
    \QCoh(X)
\]
is strong symmetric monoidal.

Every object of
\[
    \Perf(X)
\]
is dualizable by definition, so
\[
    \Perf(X)
\]
is rigid. For
\[
    P,M,N\in\Perf(X),
\]
duality gives natural equivalences
\[
    \Map_{\Perf(X)}
    \bigl(
        P,
        M^\vee\otimes N
    \bigr)
    \simeq
    \Map_{\Perf(X)}
    \bigl(
        P\otimes M,
        N
    \bigr).
\]
Thus
\[
    M^\vee\otimes N
\]
represents the internal Hom functor in
\[
    \Perf(X),
\]
and the canonical closed self-enrichment has
\[
    \underline{\Perf(X)}(M,N)
    \simeq
    M^\vee\otimes N.
\]

Finally, the preceding argument shows that the limit formed in
\[
    \CAlg(\Cat_\infty)
\]
is stable and idempotent complete and has tensor product exact separately.
A strong symmetric monoidal functor into this limit is exact if and only if
its composites with all affine evaluation functors are exact, since finite
limits and colimits are computed componentwise. Hence this object also
satisfies the universal property of the limit in
\[
    \CAlg(\Cat_\infty^{\mathrm{perf}}).
\]
This gives the asserted upgrade.
\end{proof}

\begin{proposition}[Functoriality of perfect objects]
\label{prop:perfect-stack-functoriality}
The quasi-coherent pullback functors restrict to a functor
\[
    \Perf
    \colon
    \PSh(\Aff)^{\op}
    \longrightarrow
    \CAlg(\Cat_\infty^{\mathrm{perf}}).
\]

For a morphism
\[
    p
    \colon
    X
    \longrightarrow
    Y,
\]
its value is the restriction
\[
    p^*
    \colon
    \Perf(Y)
    \longrightarrow
    \Perf(X)
\]
of quasi-coherent pullback.
\end{proposition}

\begin{proof}
The functor
\[
    p^*
    \colon
    \QCoh(Y)
    \longrightarrow
    \QCoh(X)
\]
is strong symmetric monoidal. It therefore preserves dualizable objects
and their duality data and hence restricts to
\[
    p^*
    \colon
    \Perf(Y)
    \longrightarrow
    \Perf(X).
\]

It remains to verify exactness. For every affine point
\[
    U
    \longrightarrow
    X,
\]
evaluation satisfies
\[
    \ev_U\circ p^*
    \simeq
    \ev_{p\circ U}.
\]
Zero objects, fibres, and cofibres in the affine-point limit descriptions
are computed componentwise. Therefore
\[
    p^*
\]
preserves zero objects and cofiber sequences after evaluation at every
affine point. Since affine evaluations detect equivalences, it follows
that
\[
    p^*
    \colon
    \Perf(Y)
    \longrightarrow
    \Perf(X)
\]
is exact.

The composition, unit, symmetric monoidal, and higher coherence data are
the restrictions of those of
\[
    \QCoh
    \colon
    \PSh(\Aff)^{\op}
    \longrightarrow
    \CAlg(\Cat_\infty).
\]
Thus the assignments define the stated functor.
\end{proof}

\section{Rigid self-enrichment of perfect objects}
\label{sec:perfect-self-enrichment}


\subsection{Affine rigid self-enrichment}

\begin{theorem}[Perfect self-enrichment]
\label{thm:perfect-self-enrichment}
The rigid symmetric monoidal $\infty$-category
\[
    \Perf(A)
\]
is enriched over itself. Its enriched Hom objects are
\[
    \underline{\Perf(A)}(M,N)
    :=
    M^\vee\otimes_A N.
\]
Under the equivalence of
\Ref{thm:tensor-dual-internal-hom},
this is the restriction of the closed self-enrichment of
\[
    \Mod_A(\V).
\]
\end{theorem}

\begin{proof}
By
\Ref{cor:perfect-internal-hom},
the internal Hom of two perfect modules is again perfect. Thus the
self-enrichment of
\[
    \Mod_A(\V)
\]
restricts to
\[
    \Perf(A).
\]

The enriched identity of
\[
    M
\]
is the composite
\[
    A
    \xrightarrow{\coev_M}
    M\otimes_A M^\vee
    \xrightarrow{\beta_{M,M^\vee}}
    M^\vee\otimes_A M.
\]

For
\[
    L,M,N\in\Perf(A),
\]
enriched composition is the contraction
\[
\begin{aligned}
    &(M^\vee\otimes_A N)
    \otimes_A
    (L^\vee\otimes_A M)
    \\
    &\qquad\xrightarrow{\;\simeq\;}
    L^\vee
    \otimes_A
    M^\vee
    \otimes_A
    M
    \otimes_A
    N
    \\
    &\qquad\xrightarrow{
        \id_{L^\vee}
        \otimes
        \ev_M
        \otimes
        \id_N
    }
    L^\vee\otimes_A N.
\end{aligned}
\]
Associativity and the identity laws are the restrictions of those of the
ambient closed self-enrichment. Equivalently, they follow from the
triangle identities and the coherence of the symmetric monoidal
structure.
\end{proof}


\subsection{Base change of enriched Hom objects}

\begin{proposition}[Base change of enriched Hom objects]
\label{prop:base-change-enriched-hom}
Let
\[
    A
    \longrightarrow
    B
\]
be a morphism in
\[
    \CAlg(\V),
\]
and let
\[
    M,N\in\Perf(A).
\]
Then there is a canonical equivalence
\[
\begin{aligned}
    &
    \underline{\Perf(B)}
    \bigl(
        M\otimes_A B,
        N\otimes_A B
    \bigr)
    \\
    &\qquad\simeq
    \underline{\Perf(A)}(M,N)
    \otimes_A B.
\end{aligned}
\]
\end{proposition}

\begin{proof}
Using
\Ref{prop:base-change-perfect-modules},
\[
\begin{aligned}
    &
    \underline{\Perf(B)}
    \bigl(
        M\otimes_A B,
        N\otimes_A B
    \bigr)
    \\
    &\qquad\simeq
    (M^\vee\otimes_A B)
    \otimes_B
    (N\otimes_A B)
    \\
    &\qquad\simeq
    (M^\vee\otimes_A N)
    \otimes_A B
    \\
    &\qquad=
    \underline{\Perf(A)}(M,N)
    \otimes_A B.
\end{aligned}
\]
\end{proof}


\subsection{Stackwise rigid self-enrichment}

\begin{proposition}[Self-enrichment on a relative stack]
\label{prop:perfect-stack-self-enrichment}
For every
\[
    X\in\PSh(\Aff),
\]
the rigid symmetric monoidal category
\[
    \Perf(X)
\]
is enriched over itself with
\[
    \underline{\Perf(X)}(M,N)
    :=
    M^\vee\otimes N.
\]
Its pullback to every affine point
\[
    U
    \longrightarrow
    X
\]
agrees with the affine enriched Hom object.
\end{proposition}

\begin{proof}
The construction of enriched identities and composition uses only tensor
products, evaluation, coevaluation, and the symmetry. These structures
exist in
\[
    \Perf(X)
\]
and are computed componentwise in the affine-point limit description of
\Ref{prop:perfect-limit-description}.

For every affine point
\[
    U
    \longrightarrow
    X,
\]
the evaluation functor
\[
    \Perf(X)
    \longrightarrow
    \Perf(U)
\]
is strong symmetric monoidal. It therefore preserves duals, evaluation,
coevaluation, enriched identities, and enriched composition. Hence it
carries
\[
    M^\vee\otimes N
\]
to
\[
    (M|_U)^\vee\otimes(N|_U).
\]
\end{proof}

\begin{remark}
\label{rem:perfect-restriction-after-construction}
The ambient module category is constructed before restricting to perfect
objects. The category
\[
    \Perf(A)
\]
need not possess the geometric realizations required to construct
arbitrary relative tensor products of associative modules or general bar
constructions. Such constructions are performed in the presentable
category
\[
    \Mod_A(\V)
\]
and then restricted when the resulting objects are perfect.
\end{remark}

\section{The canonical symmetric bilinear functor and coherent duality}
\label{sec:canonical-perfect-duality}

Set
\[
    \mathcal C_A
    :=
    \Perf(A).
\]


\subsection{The canonical bilinear mapping spectrum}

\begin{definition}[Canonical bilinear mapping spectrum]
\label{def:canonical-perfect-bilinear-functor}
For
\[
    M,N\in\mathcal C_A,
\]
define
\[
    \mathbb B_A(M,N)
    :=
    \map_{\mathcal C_A}
    \bigl(
        M\otimes_A N,
        A
    \bigr).
\]
\end{definition}

\begin{proposition}[Symmetry and representability]
\label{prop:canonical-bilinear-symmetry-representability}
The functor
\[
    \mathbb B_A
    \colon
    \mathcal C_A^{\op}
    \times
    \mathcal C_A^{\op}
    \longrightarrow
    \Sp
\]
is exact separately in each variable.

The braiding of
\[
    \mathcal C_A
\]
induces a coherent symmetry equivalence
\[
    s_{M,N}
    \colon
    \mathbb B_A(M,N)
    \xrightarrow{\;\simeq\;}
    \mathbb B_A(N,M).
\]

Moreover,
\[
    \mathbb B_A
\]
is right represented by dualization:
\[
    \mathbb B_A(M,N)
    \simeq
    \map_{\mathcal C_A}
    \bigl(
        M,
        N^\vee
    \bigr).
\]
\end{proposition}

\begin{proof}
Tensoring is exact separately in each variable, and mapping into
\[
    A
\]
is exact contravariantly. Hence
\[
    \mathbb B_A
\]
is exact separately in each variable.

The braiding
\[
    \beta_{N,M}
    \colon
    N\otimes_A M
    \xrightarrow{\;\simeq\;}
    M\otimes_A N
\]
induces
\[
    s_{M,N}(\varphi)
    :=
    \varphi\circ\beta_{N,M}.
\]
The full symmetric monoidal coherence of the braiding supplies
\[
    s
\]
with its coherent involutivity and all higher compatibilities. Thus
\[
    \mathbb B_A
\]
is a symmetric bilinear functor.

Finally, rigidity gives natural equivalences
\[
\begin{aligned}
    \mathbb B_A(M,N)
    &=
    \map_{\mathcal C_A}
    \bigl(
        M\otimes_A N,
        A
    \bigr)
    \\
    &\simeq
    \map_{\mathcal C_A}
    \bigl(
        M,
        N^\vee
    \bigr).
\end{aligned}
\]
\end{proof}


\subsection{Coherent perfect dualities}

\begin{definition}[Coherent perfect duality]
\label{def:coherent-perfect-duality}
Let
\[
    \mathcal C
\]
be a small stable $\infty$-category. A
\textbf{coherent perfect duality} on
\[
    \mathcal C
\]
consists of an exact equivalence
\[
    D
    \colon
    \mathcal C^{\op}
    \xrightarrow{\;\simeq\;}
    \mathcal C,
\]
a natural equivalence
\[
    \operatorname{bid}
    \colon
    \id_{\mathcal C}
    \xrightarrow{\;\simeq\;}
    D D^{\op},
\]
and coherent higher
\[
    C_2
\]
-data whose first nontrivial coherence is
\[
    D(\operatorname{bid}_X)
    \circ
    \operatorname{bid}_{D(X)}
    \simeq
    \id_{D(X)}
\]
for every
\[
    X\in\mathcal C.
\]

Here
\[
    D^{\op}
    \colon
    \mathcal C
    \longrightarrow
    \mathcal C^{\op}
\]
is the opposite functor, so that
\[
    D D^{\op}
    \colon
    \mathcal C
    \longrightarrow
    \mathcal C.
\]

Equivalently, a coherent perfect duality is a homotopy-fixed-point
structure for the involution on small stable $\infty$-categories induced
by passage to the opposite category, whose underlying equivalence is
\[
    D
    \colon
    \mathcal C^{\op}
    \xrightarrow{\;\simeq\;}
    \mathcal C.
\]
\end{definition}


\subsection{The canonical coherent duality}

\begin{proposition}[Canonical perfect duality]
\label{prop:canonical-perfect-duality}
Dualization defines an exact contravariant equivalence
\[
    D_A
    :=
    (-)^\vee
    \colon
    \Perf(A)^{\op}
    \xrightarrow{\;\simeq\;}
    \Perf(A).
\]

The coherent symmetry of
\[
    \mathbb B_A
\]
canonically promotes
\[
    D_A
\]
to a coherent perfect duality. Its biduality transformation is the
canonical natural equivalence
\[
    \operatorname{bid}_M
    \colon
    M
    \xrightarrow{\;\simeq\;}
    M^{\vee\vee}.
\]
It satisfies
\[
    D_A(\operatorname{bid}_M)
    \circ
    \operatorname{bid}_{D_A(M)}
    \simeq
    \id_{D_A(M)}
\]
together with all higher
\[
    C_2
\]
-coherences.
\end{proposition}

\begin{proof}
Under the representing equivalence of
\Ref{prop:canonical-bilinear-symmetry-representability},
the symmetry of
\[
    \mathbb B_A
\]
gives natural equivalences
\[
    \theta_{M,N}
    \colon
    \map_{\mathcal C_A}
    \bigl(
        M,
        D_A N
    \bigr)
    \xrightarrow{\;\simeq\;}
    \map_{\mathcal C_A}
    \bigl(
        N,
        D_A M
    \bigr).
\]

Explicitly, if
\[
    q
    \colon
    M
    \longrightarrow
    N^\vee,
\]
then
\[
    \theta_{M,N}(q)
    =
    q^\vee
    \circ
    \operatorname{bid}_N
    \colon
    N
    \longrightarrow
    M^\vee.
\]
Indeed, currying
\[
    q
\]
gives the pairing
\[
    M\otimes_A N
    \xrightarrow{q\otimes\id_N}
    N^\vee\otimes_A N
    \xrightarrow{\ev_N}
    A.
\]
After precomposition with the braiding
\[
    N\otimes_A M
    \xrightarrow{\beta_{N,M}}
    M\otimes_A N,
\]
currying in the
\[
    M
\]
-variable gives
\[
    q^\vee\circ\operatorname{bid}_N.
\]

Let
\[
    D_A^{\op}
    \colon
    \mathcal C_A
    \longrightarrow
    \mathcal C_A^{\op}
\]
be the opposite functor. For
\[
    M\in\mathcal C_A
\]
and
\[
    N\in\mathcal C_A^{\op},
\]
there are natural equivalences
\[
\begin{aligned}
    \map_{\mathcal C_A^{\op}}
    \bigl(
        D_A^{\op}M,
        N
    \bigr)
    &=
    \map_{\mathcal C_A}
    \bigl(
        N,
        D_A M
    \bigr)
    \\
    &\simeq
    \map_{\mathcal C_A}
    \bigl(
        M,
        D_A N
    \bigr).
\end{aligned}
\]
Thus
\[
    D_A^{\op}
    \dashv
    D_A.
\]

The unit
\[
    \eta_M
    \colon
    M
    \longrightarrow
    D_A^2M
\]
is characterized by
\[
    \theta_{M,D_AM}(\eta_M)
    \simeq
    \id_{D_AM}.
\]
The involutivity of
\[
    \theta
\]
gives
\[
    \theta_{M,D_AM}^{-1}
    \simeq
    \theta_{D_AM,M},
\]
and therefore
\[
    \eta_M
    =
    \theta_{D_AM,M}
    \bigl(
        \id_{D_AM}
    \bigr).
\]
By the explicit formula for
\[
    \theta,
\]
this is precisely
\[
    \operatorname{bid}_M
    \colon
    M
    \longrightarrow
    M^{\vee\vee}.
\]

Let
\[
    \epsilon_M
    \colon
    D_A^{\op}D_A M
    \longrightarrow
    M
\]
be the counit in
\[
    \mathcal C_A^{\op}.
\]
Its opposite is a morphism
\[
    \epsilon_M^{\op}
    \colon
    M
    \longrightarrow
    D_A^2M
\]
in
\[
    \mathcal C_A.
\]
The defining property of the counit gives
\[
    \epsilon_M^{\op}
    =
    \theta_{D_AM,M}
    \bigl(
        \id_{D_AM}
    \bigr)
    =
    \operatorname{bid}_M.
\]
Thus the unit and the opposite of the counit are both the canonical
biduality map.

Since every object of
\[
    \Perf(A)
\]
is dualizable,
\[
    \operatorname{bid}_M
\]
is an equivalence. Hence
\[
    D_A^{\op}
    \dashv
    D_A
\]
is an adjoint equivalence.

The coherent symmetry of
\[
    \mathbb B_A
\]
is a homotopy fixed point for the flip action on perfect bilinear
functors. The equivalence between perfect symmetric bilinear functors and
perfect dualities transports this entire fixed-point datum to
\[
    D_A.
\]
It therefore supplies the displayed biduality coherence and all higher
\[
    C_2
\]
-coherences. See
\cite{CalmesDottoHarpazEtAl}.

Finally,
\[
    D_A
\]
is an equivalence between stable $\infty$-categories and is therefore
exact.
\end{proof}


\subsection{Evaluation and coevaluation under biduality}

\begin{lemma}[Evaluation and coevaluation under biduality]
\label{lem:biduality-evaluation-coevaluation}
For every
\[
    M\in\Perf(A),
\]
there are canonical equivalences
\[
    \ev_{M^\vee}
    \circ
    \bigl(
        \operatorname{bid}_M
        \otimes_A
        \id_{M^\vee}
    \bigr)
    \simeq
    \ev_M
    \circ
    \beta_{M,M^\vee}
\]
as morphisms
\[
    M\otimes_A M^\vee
    \longrightarrow
    A,
\]
and
\[
    \bigl(
        \operatorname{bid}_M
        \otimes_A
        \id_{M^\vee}
    \bigr)
    \circ
    \coev_M
    \simeq
    \beta_{M^\vee,M^{\vee\vee}}
    \circ
    \coev_{M^\vee}
\]
as morphisms
\[
    A
    \longrightarrow
    M^{\vee\vee}\otimes_A M^\vee.
\]
\end{lemma}

\begin{proof}
The biduality map
\[
    \operatorname{bid}_M
    \colon
    M
    \longrightarrow
    M^{\vee\vee}
\]
is the adjunct, under
\[
    \Map_{\Perf(A)}
    \bigl(
        M,
        M^{\vee\vee}
    \bigr)
    \simeq
    \Map_{\Perf(A)}
    \bigl(
        M\otimes_A M^\vee,
        A
    \bigr),
\]
of the transposed evaluation pairing
\[
    M\otimes_A M^\vee
    \xrightarrow{\beta_{M,M^\vee}}
    M^\vee\otimes_A M
    \xrightarrow{\ev_M}
    A.
\]
The defining property of this adjunct is exactly
\[
    \ev_{M^\vee}
    \circ
    \bigl(
        \operatorname{bid}_M
        \otimes_A
        \id_{M^\vee}
    \bigr)
    \simeq
    \ev_M
    \circ
    \beta_{M,M^\vee}.
\]

We now derive the coevaluation identity. Define
\[
\begin{aligned}
    c_M
    :={}&
    \beta_{M^{\vee\vee},M^\vee}
    \circ
    \bigl(
        \operatorname{bid}_M
        \otimes_A
        \id_{M^\vee}
    \bigr)
    \circ
    \coev_M
    \\
    &\colon
    A
    \longrightarrow
    M^\vee\otimes_A M^{\vee\vee}.
\end{aligned}
\]
By naturality of the braiding,
\[
    c_M
    \simeq
    \bigl(
        \id_{M^\vee}
        \otimes_A
        \operatorname{bid}_M
    \bigr)
    \circ
    \beta_{M,M^\vee}
    \circ
    \coev_M.
\]

Consider the adjunction equivalence
\[
\begin{aligned}
    \Phi
    \colon
    \Map_{\Perf(A)}
    \bigl(
        A,
        M^\vee\otimes_A M^{\vee\vee}
    \bigr)
    \xrightarrow{\;\simeq\;}
    \Map_{\Perf(A)}(M^\vee,M^\vee)
\end{aligned}
\]
given by
\[
\begin{aligned}
    \Phi(c)
    :=
    \Bigl[
    M^\vee
    &\simeq
    A\otimes_A M^\vee
    \\
    &\xrightarrow{
        c\otimes_A\id_{M^\vee}
    }
    M^\vee
    \otimes_A
    M^{\vee\vee}
    \otimes_A
    M^\vee
    \\
    &\xrightarrow{
        \id_{M^\vee}
        \otimes_A
        \ev_{M^\vee}
    }
    M^\vee
    \Bigr].
\end{aligned}
\]

By the triangle identity,
\[
    \Phi(\coev_{M^\vee})
    \simeq
    \id_{M^\vee}.
\]
For
\[
    c_M,
\]
the corresponding composite contains the contraction
\[
    M\otimes_A M^\vee
    \xrightarrow{
        \operatorname{bid}_M
        \otimes_A
        \id_{M^\vee}
    }
    M^{\vee\vee}\otimes_A M^\vee
    \xrightarrow{\ev_{M^\vee}}
    A.
\]
By the first equivalence of the lemma, this contraction is canonically
equivalent to
\[
    M\otimes_A M^\vee
    \xrightarrow{\beta_{M,M^\vee}}
    M^\vee\otimes_A M
    \xrightarrow{\ev_M}
    A.
\]
After this replacement, the composite defining
\[
    \Phi(c_M)
\]
is the triangle composite for the dual pair
\[
    (M,M^\vee).
\]
Hence
\[
    \Phi(c_M)
    \simeq
    \id_{M^\vee}.
\]
Since
\[
    \Phi
\]
is an equivalence,
\[
    c_M
    \simeq
    \coev_{M^\vee}.
\]
Composing with
\[
    \beta_{M^\vee,M^{\vee\vee}}
\]
gives
\[
    \bigl(
        \operatorname{bid}_M
        \otimes_A
        \id_{M^\vee}
    \bigr)
    \circ
    \coev_M
    \simeq
    \beta_{M^\vee,M^{\vee\vee}}
    \circ
    \coev_{M^\vee}.
\]
\end{proof}


\section{Functoriality of the canonical duality}
\label{sec:functoriality-canonical-perfect-duality}
\label{sec:base-change-perfect-duality}

The coherent duality constructed above is functorial under exact strong
symmetric monoidal functors between rigid stable symmetric monoidal
$\infty$-categories. We record the complete comparison because it will be
used both for extension of scalars and for pullback on relative stacks.

\begin{proposition}[Duality preservation by strong symmetric monoidal functors]
\label{prop:duality-preserving-strong-monoidal}
Let
\[
    F
    \colon
    \mathcal C
    \longrightarrow
    \mathcal D
\]
be an exact strong symmetric monoidal functor between rigid stable
symmetric monoidal $\infty$-categories. For every
\[
    M\in\mathcal C,
\]
there is a canonical natural equivalence
\[
    \chi_M
    \colon
    F(M^\vee)
    \xrightarrow{\;\simeq\;}
    F(M)^\vee.
\]
The family
\[
    \chi_M
\]
is compatible with evaluation, coevaluation, composition, units, and
biduality. More precisely, for every
\[
    M\in\mathcal C,
\]
there is a canonical equivalence
\[
\begin{aligned}
    &
    (\chi_M^{-1})^\vee
    \circ
    \chi_{M^\vee}
    \circ
    F(\operatorname{bid}_M)
    \\
    &\qquad\simeq
    \operatorname{bid}_{F(M)}
\end{aligned}
\]
as morphisms
\[
    F(M)
    \longrightarrow
    F(M)^{\vee\vee}.
\]
Equivalently,
\[
    \chi_M^\vee
    \circ
    \operatorname{bid}_{F(M)}
    \simeq
    \chi_{M^\vee}
    \circ
    F(\operatorname{bid}_M).
\]
Together with the induced higher coherences, these comparisons make
\[
    F
\]
a duality-preserving exact strong symmetric monoidal functor.
\end{proposition}

\begin{proof}
If
\[
    M^\vee
\]
is a dual of
\[
    M,
\]
then
\[
    F(M^\vee)
\]
is a dual of
\[
    F(M).
\]
Its evaluation is
\[
\begin{aligned}
    F(M^\vee)\otimes F(M)
    &\simeq
    F(M^\vee\otimes M)
    \\
    &\xrightarrow{F(\ev_M)}
    F(\mathbf 1_{\mathcal C})
    \simeq
    \mathbf 1_{\mathcal D},
\end{aligned}
\]
and its coevaluation is
\[
\begin{aligned}
    \mathbf 1_{\mathcal D}
    &\simeq
    F(\mathbf 1_{\mathcal C})
    \\
    &\xrightarrow{F(\coev_M)}
    F(M\otimes M^\vee)
    \\
    &\simeq
    F(M)\otimes F(M^\vee).
\end{aligned}
\]

The canonical object
\[
    F(M)^\vee
\]
is another dual of
\[
    F(M).
\]
The contractible uniqueness of duals therefore supplies a canonical
equivalence
\[
    \chi_M
    \colon
    F(M^\vee)
    \xrightarrow{\;\simeq\;}
    F(M)^\vee
\]
compatible with evaluation and coevaluation.

The construction is natural in
\[
    M.
\]
Indeed, for a morphism
\[
    u
    \colon
    M
    \longrightarrow
    N,
\]
both composites
\[
    F(N^\vee)
    \xrightarrow{F(u^\vee)}
    F(M^\vee)
    \xrightarrow{\chi_M}
    F(M)^\vee
\]
and
\[
    F(N^\vee)
    \xrightarrow{\chi_N}
    F(N)^\vee
    \xrightarrow{F(u)^\vee}
    F(M)^\vee
\]
are the morphism of duals mate to
\[
    F(u).
\]
Hence there is a canonical equivalence
\[
    \chi_M
    \circ
    F(u^\vee)
    \simeq
    F(u)^\vee
    \circ
    \chi_N.
\]

The comparisons also preserve biduality. Both
\[
    (\chi_M^{-1})^\vee
    \circ
    \chi_{M^\vee}
    \circ
    F(\operatorname{bid}_M)
\]
and
\[
    \operatorname{bid}_{F(M)}
\]
are the canonical morphism from
\[
    F(M)
\]
to the double dual determined by the transported dual pair. Contractible
uniqueness of the complete dual-pair datum identifies them, giving
\[
\begin{aligned}
    &
    (\chi_M^{-1})^\vee
    \circ
    \chi_{M^\vee}
    \circ
    F(\operatorname{bid}_M)
    \\
    &\qquad\simeq
    \operatorname{bid}_{F(M)}.
\end{aligned}
\]
After composing with
\[
    \chi_M^\vee,
\]
this is equivalently
\[
    \chi_M^\vee
    \circ
    \operatorname{bid}_{F(M)}
    \simeq
    \chi_{M^\vee}
    \circ
    F(\operatorname{bid}_M).
\]

The comparisons and their composition, unit, and higher coherences are
induced functorially from the walking dual pair. Thus
\[
    F
\]
is a duality-preserving exact strong symmetric monoidal functor.
\end{proof}

\begin{corollary}[Base change of perfect duality]
\label{cor:base-change-perfect-duality}
Let
\[
    A
    \longrightarrow
    A'
\]
be a morphism in
\[
    \CAlg(\V).
\]
For every
\[
    M\in\Perf(A),
\]
there is a canonical equivalence
\[
    (M\otimes_A A')^\vee
    \simeq
    M^\vee\otimes_A A'.
\]
The resulting comparisons make extension of scalars
\[
    -\otimes_A A'
    \colon
    \Perf(A)
    \longrightarrow
    \Perf(A')
\]
a duality-preserving exact strong symmetric monoidal functor.
\end{corollary}

\begin{proof}
By
\Ref{prop:indexed-perfect-category},
extension of scalars is exact and strong symmetric monoidal. Apply
\Ref{prop:duality-preserving-strong-monoidal}.
\end{proof}

\begin{corollary}[Pullback compatibility of perfect duality]
\label{cor:pullback-perfect-duality}
Let
\[
    p
    \colon
    X
    \longrightarrow
    Y
\]
be a morphism in
\[
    \PSh(\Aff).
\]
For every
\[
    M\in\Perf(Y),
\]
there is a canonical natural equivalence
\[
    \chi_{p,M}
    \colon
    p^*(M^\vee)
    \xrightarrow{\;\simeq\;}
    \bigl(p^*M\bigr)^\vee.
\]
These comparisons make
\[
    p^*
    \colon
    \Perf(Y)
    \longrightarrow
    \Perf(X)
\]
a duality-preserving exact strong symmetric monoidal functor, coherently
compatible with identities and composition of morphisms of presheaves.
\end{corollary}

\begin{proof}
By
\Ref{prop:perfect-stack-functoriality},
quasi-coherent pullback restricts to an exact strong symmetric monoidal
functor on perfect objects. Apply
\Ref{prop:duality-preserving-strong-monoidal}.
The identity and composition coherences are those induced functorially from
the walking dual pair together with the corresponding coherences of
quasi-coherent pullback.
\end{proof}

\section{Enriched duality, cocomposition, and trace}
\label{sec:canonical-perfect-duality-trace}


\subsection{Duality symmetry of enriched Hom objects}

\begin{theorem}[Duality symmetry of enriched Hom objects]
\label{thm:duality-enriched-hom}
For
\[
    M,N\in\Perf(A),
\]
there is a canonical natural equivalence
\[
    \delta_{M,N}
    \colon
    \underline{\Perf(A)}(M,N)^\vee
    \xrightarrow{\;\simeq\;}
    \underline{\Perf(A)}(N,M).
\]
Explicitly,
\[
    (M^\vee\otimes_A N)^\vee
    \simeq
    N^\vee\otimes_A M.
\]
\end{theorem}

\begin{proof}
There are canonical equivalences
\[
\begin{aligned}
    (M^\vee\otimes_A N)^\vee
    &\simeq
    N^\vee\otimes_A M^{\vee\vee}
    \\
    &\xrightarrow{
        \id_{N^\vee}
        \otimes
        \operatorname{bid}_M^{-1}
    }
    N^\vee\otimes_A M.
\end{aligned}
\]
The construction is natural in
\[
    M
    \qquad\text{and}\qquad
    N.
\]
\end{proof}

The equivalence
\[
    \delta_{M,N}
\]
identifies the dual of the enriched Hom object with the reverse enriched
Hom object. The enriched action of the duality functor itself is a
different map.


\subsection{Enriched dualization and reversed composition}

\begin{proposition}[Enriched dualization and reversed composition]
\label{prop:coherence-enriched-duality}
For
\[
    M,N\in\Perf(A),
\]
define
\[
    d_{M,N}
    \colon
    \underline{\Perf(A)}(M,N)
    \longrightarrow
    \underline{\Perf(A)}(N^\vee,M^\vee)
\]
by
\[
\begin{aligned}
    M^\vee\otimes_A N
    &\xrightarrow{\beta_{M^\vee,N}}
    N\otimes_A M^\vee
    \\
    &\xrightarrow{
        \operatorname{bid}_N
        \otimes
        \id_{M^\vee}
    }
    N^{\vee\vee}\otimes_A M^\vee.
\end{aligned}
\]
Then:
\begin{enumerate}[label=(\roman*)]
    \item
    \[
        d_{M,N}
    \]
    is an equivalence natural in
    \[
        M
        \qquad\text{and}\qquad
        N;
    \]

    \item
    \[
        d_{M,M}
    \]
    carries the enriched identity of
    \[
        M
    \]
    to the enriched identity of
    \[
        M^\vee;
    \]

    \item for
    \[
        L,M,N\in\Perf(A),
    \]
    the diagram
    \[
    \begin{tikzcd}[column sep=huge]
    \underline{\Perf(A)}(M,N)
    \otimes_A
    \underline{\Perf(A)}(L,M)
        \arrow[r,"\mathrm{comp}"]
        \arrow[d,"d_{M,N}\otimes d_{L,M}"']
    &
    \underline{\Perf(A)}(L,N)
        \arrow[d,"d_{L,N}"]
    \\
    \underline{\Perf(A)}(N^\vee,M^\vee)
    \otimes_A
    \underline{\Perf(A)}(M^\vee,L^\vee)
        \arrow[d,"\beta"']
    &
    \underline{\Perf(A)}(N^\vee,L^\vee)
    \\
    \underline{\Perf(A)}(M^\vee,L^\vee)
    \otimes_A
    \underline{\Perf(A)}(N^\vee,M^\vee)
        \arrow[r,"\mathrm{comp}"']
    &
    \underline{\Perf(A)}(N^\vee,L^\vee)
        \arrow[u,equal]
    \end{tikzcd}
    \]
    commutes up to its canonical coherent homotopy.
\end{enumerate}
\end{proposition}

\begin{proof}
The first assertion follows because the symmetry and biduality maps are
equivalences.

Let
\[
    u_M
    \colon
    A
    \longrightarrow
    M^\vee\otimes_A M
\]
be the enriched identity. By definition,
\[
    u_M
    =
    \beta_{M,M^\vee}
    \circ
    \coev_M.
\]
Therefore
\[
\begin{aligned}
    d_{M,M}\circ u_M
    &=
    \bigl(
        \operatorname{bid}_M
        \otimes
        \id_{M^\vee}
    \bigr)
    \circ
    \beta_{M^\vee,M}
    \circ
    \beta_{M,M^\vee}
    \circ
    \coev_M
    \\
    &\simeq
    \bigl(
        \operatorname{bid}_M
        \otimes
        \id_{M^\vee}
    \bigr)
    \circ
    \coev_M
    \\
    &\simeq
    \beta_{M^\vee,M^{\vee\vee}}
    \circ
    \coev_{M^\vee}
    \\
    &=
    u_{M^\vee},
\end{aligned}
\]
where the penultimate equivalence is
\Ref{lem:biduality-evaluation-coevaluation}.

It remains to verify compatibility with composition. Both paths are
morphisms
\[
    M^\vee
    \otimes_A
    N
    \otimes_A
    L^\vee
    \otimes_A
    M
    \longrightarrow
    N^{\vee\vee}
    \otimes_A
    L^\vee.
\]

The upper path is canonically equivalent to
\[
\begin{aligned}
    M^\vee
    \otimes_A
    N
    \otimes_A
    L^\vee
    \otimes_A
    M
    &\xrightarrow{\;\simeq\;}
    L^\vee
    \otimes_A
    M^\vee
    \otimes_A
    M
    \otimes_A
    N
    \\
    &\xrightarrow{
        \id_{L^\vee}
        \otimes
        \ev_M
        \otimes
        \id_N
    }
    L^\vee\otimes_A N
    \\
    &\xrightarrow{\beta}
    N\otimes_A L^\vee
    \\
    &\xrightarrow{
        \operatorname{bid}_N
        \otimes
        \id_{L^\vee}
    }
    N^{\vee\vee}\otimes_A L^\vee.
\end{aligned}
\]

After applying
\[
    d_{M,N}
    \qquad\text{and}\qquad
    d_{L,M},
\]
braiding the two enriched Hom objects, and composing, the lower path is
canonically equivalent to
\[
\begin{aligned}
    M^\vee
    \otimes_A
    N
    \otimes_A
    L^\vee
    \otimes_A
    M
    &\xrightarrow{\;\simeq\;}
    N^{\vee\vee}
    \otimes_A
    M^{\vee\vee}
    \otimes_A
    M^\vee
    \otimes_A
    L^\vee
    \\
    &\xrightarrow{
        \id_{N^{\vee\vee}}
        \otimes
        \ev_{M^\vee}
        \otimes
        \id_{L^\vee}
    }
    N^{\vee\vee}\otimes_A L^\vee.
\end{aligned}
\]
The middle tensor factors arise from
\[
    M\otimes_A M^\vee
    \xrightarrow{
        \operatorname{bid}_M
        \otimes
        \id_{M^\vee}
    }
    M^{\vee\vee}\otimes_A M^\vee.
\]
By
\Ref{lem:biduality-evaluation-coevaluation},
\[
    \ev_{M^\vee}
    \circ
    \bigl(
        \operatorname{bid}_M
        \otimes
        \id_{M^\vee}
    \bigr)
    \simeq
    \ev_M
    \circ
    \beta_{M,M^\vee}.
\]
Substituting this equivalence transforms the lower composite into the
upper composite. The remaining identifications are the canonical
associativity and symmetry coherences.
\end{proof}

\begin{remark}
\label{rem:enriched-mate-not-parameter-dual}
The maps
\[
    d_{M,N}
\]
assemble into an enriched contravariant equivalence. On morphisms
represented by the tensor unit they induce ordinary dualization. More
generally, they induce the corresponding enriched mate operation directly
on parameterized enriched morphisms; the parameter object itself is not
dualized.
\end{remark}


\subsection{Dual cocomposition}

\begin{proposition}[Dual cocomposition]
\label{prop:dual-cocomposition}
For
\[
    L,M,N\in\Perf(A),
\]
dualizing enriched composition and applying the equivalences
\[
    \delta
\]
gives a natural cocomposition morphism
\[
\begin{aligned}
    \Delta_{L,M,N}
    \colon
    \underline{\Perf(A)}(N,L)
    \longrightarrow{}&
    \underline{\Perf(A)}(M,L)
    \otimes_A
    \underline{\Perf(A)}(N,M).
\end{aligned}
\]
Explicitly, it is the composite
\[
\begin{aligned}
    N^\vee\otimes_A L
    &\xrightarrow{
        \id_{N^\vee\otimes_A L}
        \otimes
        \coev_M
    }
    N^\vee
    \otimes_A
    L
    \otimes_A
    M
    \otimes_A
    M^\vee
    \\
    &\xrightarrow{\;\simeq\;}
    (M^\vee\otimes_A L)
    \otimes_A
    (N^\vee\otimes_A M).
\end{aligned}
\]

Under the canonical equivalences
\[
    A^\vee
    \simeq
    A
\]
and
\[
    \delta_{M,M}
    \colon
    \underline{\Perf(A)}(M,M)^\vee
    \xrightarrow{\;\simeq\;}
    \underline{\Perf(A)}(M,M),
\]
the dual of the enriched identity is the enriched counit associated to
\[
    M,
\]
namely the evaluation morphism
\[
    \epsilon_M
    :=
    \ev_M
    \colon
    \underline{\Perf(A)}(M,M)
    \longrightarrow
    A.
\]

The morphisms
\[
    \Delta
\]
are coherently coassociative, and the morphisms
\[
    \epsilon_M
\]
satisfy the two counit identities.
\end{proposition}

\begin{proof}
Enriched composition is
\[
\begin{aligned}
    \underline{\Perf(A)}(M,N)
    \otimes_A
    \underline{\Perf(A)}(L,M)
    \longrightarrow
    \underline{\Perf(A)}(L,N).
\end{aligned}
\]
Dualizing gives
\[
\begin{aligned}
    \underline{\Perf(A)}(L,N)^\vee
    \longrightarrow
    \underline{\Perf(A)}(L,M)^\vee
    \otimes_A
    \underline{\Perf(A)}(M,N)^\vee.
\end{aligned}
\]
Applying
\Ref{thm:duality-enriched-hom}
identifies this with
\[
\begin{aligned}
    \underline{\Perf(A)}(N,L)
    \longrightarrow
    \underline{\Perf(A)}(M,L)
    \otimes_A
    \underline{\Perf(A)}(N,M).
\end{aligned}
\]
Unwinding the dual of the evaluation contraction gives the displayed
coevaluation formula.

The enriched identity is
\[
    u_M
    \colon
    A
    \longrightarrow
    M^\vee\otimes_A M.
\]
Its dual has type
\[
    u_M^\vee
    \colon
    (M^\vee\otimes_A M)^\vee
    \longrightarrow
    A^\vee.
\]
Under
\[
    \delta_{M,M}
\]
and
\[
    A^\vee\simeq A,
\]
this is the evaluation morphism
\[
    \ev_M
    \colon
    M^\vee\otimes_A M
    \longrightarrow
    A.
\]

Coassociativity and the counit identities follow by dualizing the
associativity and unitality coherences of enriched composition.
\end{proof}


\subsection{Perfect trace pairing and cyclicity}

\begin{proposition}[Perfect trace pairing and cyclicity]
\label{prop:perfect-trace-pairing}
For
\[
    M,N\in\Perf(A),
\]
define
\[
    \langle-,-\rangle_{M,N}
    \colon
    \underline{\Perf(A)}(N,M)
    \otimes_A
    \underline{\Perf(A)}(M,N)
    \longrightarrow
    A
\]
as the composite
\[
\begin{aligned}
    &
    \underline{\Perf(A)}(N,M)
    \otimes_A
    \underline{\Perf(A)}(M,N)
    \\
    &\xrightarrow{\beta}
    \underline{\Perf(A)}(M,N)
    \otimes_A
    \underline{\Perf(A)}(N,M)
    \\
    &\xrightarrow{\mathrm{comp}}
    \underline{\Perf(A)}(N,N)
    \\
    &\xrightarrow{\ev_N}
    A.
\end{aligned}
\]

The morphism
\[
    \operatorname{tr}_N
    :=
    \ev_N
    \colon
    \underline{\Perf(A)}(N,N)
    =
    N^\vee\otimes_A N
    \longrightarrow
    A
\]
is the trace functional on the enriched endomorphism object. After
precomposition with a morphism
\[
    A
    \longrightarrow
    N^\vee\otimes_A N
\]
representing an endomorphism of
\[
    N,
\]
it gives the categorical trace of that endomorphism.

The pairing
\[
    \langle-,-\rangle_{M,N}
\]
is perfect. Its adjoint is the inverse of
\[
    \delta_{M,N}
    \colon
    \underline{\Perf(A)}(M,N)^\vee
    \xrightarrow{\;\simeq\;}
    \underline{\Perf(A)}(N,M).
\]

For
\[
    L,M,N\in\Perf(A),
\]
the three contraction morphisms
\[
\begin{aligned}
    &
    \underline{\Perf(A)}(N,L)
    \otimes_A
    \underline{\Perf(A)}(M,N)
    \otimes_A
    \underline{\Perf(A)}(L,M)
    \longrightarrow
    A
\end{aligned}
\]
obtained by cyclically choosing the enriched endomorphism object on which
the trace is evaluated are canonically equivalent.

Consequently, for generalized elements
\[
    x
    \colon
    M
    \longrightarrow
    N,
\]
\[
    y
    \colon
    L
    \longrightarrow
    M,
\]
\[
    z
    \colon
    N
    \longrightarrow
    L,
\]
there are canonical equivalences
\[
    \langle z,x\circ y\rangle_{L,N}
    \simeq
    \langle x,y\circ z\rangle_{N,M}
    \simeq
    \langle y,z\circ x\rangle_{M,L}.
\]
Equivalently,
\[
    \operatorname{tr}_N
    \bigl(
        x\circ y\circ z
    \bigr)
    \simeq
    \operatorname{tr}_M
    \bigl(
        y\circ z\circ x
    \bigr)
    \simeq
    \operatorname{tr}_L
    \bigl(
        z\circ x\circ y
    \bigr).
\]

In particular, the pairing is canonically equivalent to composition in
the opposite order followed by evaluation on
\[
    M.
\]
More explicitly, the diagram
\[
\begin{tikzcd}[column sep=huge]
    \underline{\Perf(A)}(N,M)
    \otimes_A
    \underline{\Perf(A)}(M,N)
        \arrow[r, "{\langle-,-\rangle_{M,N}}"]
        \arrow[d, "\mathrm{comp}"']
    &
    A
    \\
    \underline{\Perf(A)}(M,M)
        \arrow[ur, "\ev_M"']
\end{tikzcd}
\]
commutes up to its canonical coherent homotopy.
\end{proposition}

\begin{proof}
Using
\[
    \underline{\Perf(A)}(N,M)
    =
    N^\vee\otimes_A M
\]
and
\[
    \underline{\Perf(A)}(M,N)
    =
    M^\vee\otimes_A N,
\]
the pairing is the composite
\[
\begin{aligned}
    &(N^\vee\otimes_A M)
    \otimes_A
    (M^\vee\otimes_A N)
    \\
    &\qquad\xrightarrow{\;\simeq\;}
    N^\vee
    \otimes_A
    M
    \otimes_A
    M^\vee
    \otimes_A
    N
    \\
    &\qquad\xrightarrow{
        \id_{N^\vee}
        \otimes
        \ev_M
        \otimes
        \id_N
    }
    N^\vee\otimes_A N
    \\
    &\qquad\xrightarrow{\ev_N}
    A.
\end{aligned}
\]

Its adjoint is the morphism
\[
    N^\vee\otimes_A M
    \longrightarrow
    (M^\vee\otimes_A N)^\vee.
\]
Under the canonical equivalence
\[
\begin{aligned}
    (M^\vee\otimes_A N)^\vee
    &\simeq
    N^\vee\otimes_A M^{\vee\vee}
    \\
    &\xrightarrow{
        \id_{N^\vee}
        \otimes
        \operatorname{bid}_M^{-1}
    }
    N^\vee\otimes_A M,
\end{aligned}
\]
this adjoint is the identity. Hence it is the inverse of
\[
    \delta_{M,N}
    \colon
    (M^\vee\otimes_A N)^\vee
    \xrightarrow{\;\simeq\;}
    N^\vee\otimes_A M.
\]
Therefore the pairing is perfect.

We next prove cyclicity. The tensor product of the three enriched Hom
objects is canonically equivalent to
\[
\begin{aligned}
    &
    (N^\vee\otimes_A L)
    \otimes_A
    (M^\vee\otimes_A N)
    \otimes_A
    (L^\vee\otimes_A M)
    \\
    &\qquad\simeq
    N^\vee
    \otimes_A
    L
    \otimes_A
    M^\vee
    \otimes_A
    N
    \otimes_A
    L^\vee
    \otimes_A
    M.
\end{aligned}
\]

There is a canonical total contraction
\[
\begin{aligned}
    &
    N^\vee
    \otimes_A
    L
    \otimes_A
    M^\vee
    \otimes_A
    N
    \otimes_A
    L^\vee
    \otimes_A
    M
    \\
    &\qquad\longrightarrow
    A
\end{aligned}
\]
obtained by using the symmetry to place the three dual pairs next to one
another and then applying
\[
    \ev_N,
    \qquad
    \ev_M,
    \qquad
    \ev_L.
\]

If one first composes
\[
    L
    \xrightarrow{y}
    M
    \xrightarrow{x}
    N
\]
and then pairs the result with
\[
    z
    \colon
    N
    \longrightarrow
    L,
\]
the resulting map is the total contraction closed at
\[
    N.
\]
It therefore represents
\[
    \operatorname{tr}_N
    \bigl(
        x\circ y\circ z
    \bigr).
\]

If one first composes
\[
    N
    \xrightarrow{z}
    L
    \xrightarrow{y}
    M
\]
and then pairs the result with
\[
    x
    \colon
    M
    \longrightarrow
    N,
\]
the resulting map is the same total contraction closed at
\[
    M.
\]
It therefore represents
\[
    \operatorname{tr}_M
    \bigl(
        y\circ z\circ x
    \bigr).
\]

If one first composes
\[
    M
    \xrightarrow{x}
    N
    \xrightarrow{z}
    L
\]
and then pairs the result with
\[
    y
    \colon
    L
    \longrightarrow
    M,
\]
the resulting map is the same total contraction closed at
\[
    L.
\]
It therefore represents
\[
    \operatorname{tr}_L
    \bigl(
        z\circ x\circ y
    \bigr).
\]

The three composites differ only by coherent reassociation and
permutation of tensor factors before the same evaluation contractions are
performed. The coherence theorem for symmetric monoidal categories with
duals therefore supplies canonical equivalences
\[
    \operatorname{tr}_N
    \bigl(
        x\circ y\circ z
    \bigr)
    \simeq
    \operatorname{tr}_M
    \bigl(
        y\circ z\circ x
    \bigr)
    \simeq
    \operatorname{tr}_L
    \bigl(
        z\circ x\circ y
    \bigr).
\]
All Koszul signs on homotopy groups are already encoded by the symmetry
morphisms of the stable symmetric monoidal category.

Finally, the two-factor instance of the same contraction identifies
\[
\begin{aligned}
    &
    \underline{\Perf(A)}(N,M)
    \otimes_A
    \underline{\Perf(A)}(M,N)
    \\
    &\xrightarrow{\beta}
    \underline{\Perf(A)}(M,N)
    \otimes_A
    \underline{\Perf(A)}(N,M)
    \\
    &\xrightarrow{\mathrm{comp}}
    \underline{\Perf(A)}(N,N)
    \xrightarrow{\ev_N}
    A
\end{aligned}
\]
with
\[
\begin{aligned}
    &
    \underline{\Perf(A)}(N,M)
    \otimes_A
    \underline{\Perf(A)}(M,N)
    \\
    &\xrightarrow{\mathrm{comp}}
    \underline{\Perf(A)}(M,M)
    \xrightarrow{\ev_M}
    A.
\end{aligned}
\]
This proves the final assertion.
\end{proof}



The rigid categories of perfect objects now carry their canonical coherent
dualities, enriched anti-equivalences, cocomposition maps, and perfect
cyclic trace pairings. The next chapter studies coherent fixed points of
these dualities and identifies them with the Poincar\'e objects of canonical
homotopy-symmetric quadratic functors.
